\chapter[(Джанака)]{Слава самореалізації (Джанака)}

[4.1] Ось: стійкого знавця Атмана, що грається грою насолод, годі рівняти з дурнями, які тягнуть ярмо сансари.

[4.2] Те становище, якого жадають усі боги, від Шакри почавши, в убогості своїй, --- диво: йог, що вже стоїть у ньому, не відчуває захвату.

[4.3] У того, хто це знає, дотик заслуги й гріха не виникає всередині --- як простір не єднається з димом, хоч той і видимий.

[4.4] Коли великий духом пізнав, що весь цей світ --- Атман, --- хто здатен йому заборонити те, що відбувається саме собою?

[4.5] У чотиривидному роді істот, від Брахми до травинки, --- лише в знаючого є спромога відкинути бажання й небажання.

[4.6] Хтось знає Атмана --- недвоїстого Владику світу. Що він знає, те й чинить; страх не торкається його ніде.
